% !PRINTED 157
\weberpage{157}
This remark gives the rule according to which the transposed substitutions are composed, which is expressed in the theorem:

\par
8.\enspace\gesperrt{If $A_1$, $B_1$ are the transposes of $A$, $B$, then $A_1 B_1$ is the transpose of $BA$. In symbols:}

\weq{24}{(BA)_1 = A_1 B_1.}

If we multiply the substitution formulae (1) by an undetermined factor $\eta_k$ and then take the sum with respect to $k$, it follows that:

\weq{25}{\sum y_k \eta_k = \sum^i \sum^k x_i\, a_i^{(k)}\, \eta_k,}

or, if we

\weq{26}{\sum^k a_i^{(k)} \eta_k = \xi_i}

set,

\weq{27}{\sum y_k \eta_k = \sum x_i \xi_i.}

If, then, we represent the substitutions (1) by (4), the formulae (26) are the expression for the substitution

\[
(\xi) = A_1(\eta),
\]

and we can therefore also state the theorem:

\par
9.\enspace\gesperrt{If $A$ and $A_1$ are transposes of one another of dimension $n$, and if $x$, $y$, $\xi$, $\eta$ are four systems of variables which are related to one another by the substitutions}

\weq{28}{(y) = A(x), \quad (\xi) = A_1(\eta)}

\gesperrt{then the identity holds}

\weq{29}{y_1\eta_1 + y_2\eta_2 + \cdots y_n\eta_n = x_1\xi_1 + x_2\xi_2 + \cdots x_n\xi_n.}

If two sequences of variables

\[
\begin{gathered} x_1, x_2, \ldots x_n \\ \xi_1, \xi_2, \ldots \xi_n \end{gathered}
\]

are transformed simultaneously by the substitutions (28) into two new sequences

\[
\begin{gathered} y_1, y_2, \ldots y_n \\ \eta_1, \eta_2, \ldots \eta_n \end{gathered}
\]

then the two sequences of variables are called \gesperrt{contragredient} to one another, and the formulae (28) represent two mutually \gesperrt{contragredient transformations}.
